% Options for packages loaded elsewhere
% Options for packages loaded elsewhere
\PassOptionsToPackage{unicode}{hyperref}
\PassOptionsToPackage{hyphens}{url}
\PassOptionsToPackage{dvipsnames,svgnames,x11names}{xcolor}
%
\documentclass[
  letterpaper,
  DIV=11,
  numbers=noendperiod]{scrartcl}
\usepackage{xcolor}
\usepackage{amsmath,amssymb}
\setcounter{secnumdepth}{5}
\usepackage{iftex}
\ifPDFTeX
  \usepackage[T1]{fontenc}
  \usepackage[utf8]{inputenc}
  \usepackage{textcomp} % provide euro and other symbols
\else % if luatex or xetex
  \usepackage{unicode-math} % this also loads fontspec
  \defaultfontfeatures{Scale=MatchLowercase}
  \defaultfontfeatures[\rmfamily]{Ligatures=TeX,Scale=1}
\fi
\usepackage{lmodern}
\ifPDFTeX\else
  % xetex/luatex font selection
\fi
% Use upquote if available, for straight quotes in verbatim environments
\IfFileExists{upquote.sty}{\usepackage{upquote}}{}
\IfFileExists{microtype.sty}{% use microtype if available
  \usepackage[]{microtype}
  \UseMicrotypeSet[protrusion]{basicmath} % disable protrusion for tt fonts
}{}
\makeatletter
\@ifundefined{KOMAClassName}{% if non-KOMA class
  \IfFileExists{parskip.sty}{%
    \usepackage{parskip}
  }{% else
    \setlength{\parindent}{0pt}
    \setlength{\parskip}{6pt plus 2pt minus 1pt}}
}{% if KOMA class
  \KOMAoptions{parskip=half}}
\makeatother
% Make \paragraph and \subparagraph free-standing
\makeatletter
\ifx\paragraph\undefined\else
  \let\oldparagraph\paragraph
  \renewcommand{\paragraph}{
    \@ifstar
      \xxxParagraphStar
      \xxxParagraphNoStar
  }
  \newcommand{\xxxParagraphStar}[1]{\oldparagraph*{#1}\mbox{}}
  \newcommand{\xxxParagraphNoStar}[1]{\oldparagraph{#1}\mbox{}}
\fi
\ifx\subparagraph\undefined\else
  \let\oldsubparagraph\subparagraph
  \renewcommand{\subparagraph}{
    \@ifstar
      \xxxSubParagraphStar
      \xxxSubParagraphNoStar
  }
  \newcommand{\xxxSubParagraphStar}[1]{\oldsubparagraph*{#1}\mbox{}}
  \newcommand{\xxxSubParagraphNoStar}[1]{\oldsubparagraph{#1}\mbox{}}
\fi
\makeatother


\usepackage{longtable,booktabs,array}
\newcounter{none} % for unnumbered tables
\usepackage{calc} % for calculating minipage widths
% Correct order of tables after \paragraph or \subparagraph
\usepackage{etoolbox}
\makeatletter
\patchcmd\longtable{\par}{\if@noskipsec\mbox{}\fi\par}{}{}
\makeatother
% Allow footnotes in longtable head/foot
\IfFileExists{footnotehyper.sty}{\usepackage{footnotehyper}}{\usepackage{footnote}}
\makesavenoteenv{longtable}
\usepackage{graphicx}
\makeatletter
\newsavebox\pandoc@box
\newcommand*\pandocbounded[1]{% scales image to fit in text height/width
  \sbox\pandoc@box{#1}%
  \Gscale@div\@tempa{\textheight}{\dimexpr\ht\pandoc@box+\dp\pandoc@box\relax}%
  \Gscale@div\@tempb{\linewidth}{\wd\pandoc@box}%
  \ifdim\@tempb\p@<\@tempa\p@\let\@tempa\@tempb\fi% select the smaller of both
  \ifdim\@tempa\p@<\p@\scalebox{\@tempa}{\usebox\pandoc@box}%
  \else\usebox{\pandoc@box}%
  \fi%
}
% Set default figure placement to htbp
\def\fps@figure{htbp}
\makeatother


% definitions for citeproc citations
\NewDocumentCommand\citeproctext{}{}
\NewDocumentCommand\citeproc{mm}{%
  \begingroup\def\citeproctext{#2}\cite{#1}\endgroup}
\makeatletter
 % allow citations to break across lines
 \let\@cite@ofmt\@firstofone
 % avoid brackets around text for \cite:
 \def\@biblabel#1{}
 \def\@cite#1#2{{#1\if@tempswa , #2\fi}}
\makeatother
\newlength{\cslhangindent}
\setlength{\cslhangindent}{1.5em}
\newlength{\csllabelwidth}
\setlength{\csllabelwidth}{3em}
\newenvironment{CSLReferences}[2] % #1 hanging-indent, #2 entry-spacing
 {\begin{list}{}{%
  \setlength{\itemindent}{0pt}
  \setlength{\leftmargin}{0pt}
  \setlength{\parsep}{0pt}
  % turn on hanging indent if param 1 is 1
  \ifodd #1
   \setlength{\leftmargin}{\cslhangindent}
   \setlength{\itemindent}{-1\cslhangindent}
  \fi
  % set entry spacing
  \setlength{\itemsep}{#2\baselineskip}}}
 {\end{list}}
\usepackage{calc}
\newcommand{\CSLBlock}[1]{\hfill\break\parbox[t]{\linewidth}{\strut\ignorespaces#1\strut}}
\newcommand{\CSLLeftMargin}[1]{\parbox[t]{\csllabelwidth}{\strut#1\strut}}
\newcommand{\CSLRightInline}[1]{\parbox[t]{\linewidth - \csllabelwidth}{\strut#1\strut}}
\newcommand{\CSLIndent}[1]{\hspace{\cslhangindent}#1}



\setlength{\emergencystretch}{3em} % prevent overfull lines

\providecommand{\tightlist}{%
  \setlength{\itemsep}{0pt}\setlength{\parskip}{0pt}}



 


\KOMAoption{captions}{tableheading}
\makeatletter
\@ifpackageloaded{caption}{}{\usepackage{caption}}
\AtBeginDocument{%
\ifdefined\contentsname
  \renewcommand*\contentsname{Table of contents}
\else
  \newcommand\contentsname{Table of contents}
\fi
\ifdefined\listfigurename
  \renewcommand*\listfigurename{List of Figures}
\else
  \newcommand\listfigurename{List of Figures}
\fi
\ifdefined\listtablename
  \renewcommand*\listtablename{List of Tables}
\else
  \newcommand\listtablename{List of Tables}
\fi
\ifdefined\figurename
  \renewcommand*\figurename{Figure}
\else
  \newcommand\figurename{Figure}
\fi
\ifdefined\tablename
  \renewcommand*\tablename{Table}
\else
  \newcommand\tablename{Table}
\fi
}
\@ifpackageloaded{float}{}{\usepackage{float}}
\floatstyle{ruled}
\@ifundefined{c@chapter}{\newfloat{codelisting}{h}{lop}}{\newfloat{codelisting}{h}{lop}[chapter]}
\floatname{codelisting}{Listing}
\newcommand*\listoflistings{\listof{codelisting}{List of Listings}}
\makeatother
\makeatletter
\makeatother
\makeatletter
\@ifpackageloaded{caption}{}{\usepackage{caption}}
\@ifpackageloaded{subcaption}{}{\usepackage{subcaption}}
\makeatother
\usepackage{bookmark}
\IfFileExists{xurl.sty}{\usepackage{xurl}}{} % add URL line breaks if available
\urlstyle{same}
\hypersetup{
  pdftitle={Data Acquisition and Ingredient Extraction: Building a Vocabulary of What India's Packaged Food Labels Actually Say},
  pdfauthor={Lalitha A R},
  colorlinks=true,
  linkcolor={blue},
  filecolor={Maroon},
  citecolor={Blue},
  urlcolor={Blue},
  pdfcreator={LaTeX via pandoc}}


\title{Data Acquisition and Ingredient Extraction: Building a Vocabulary
of What India's Packaged Food Labels Actually Say}
\author{Lalitha A R}
\date{2026-03-01}
\begin{document}
\maketitle

\renewcommand*\contentsname{Table of contents}
{
\setcounter{tocdepth}{3}
\tableofcontents
}

\subsection*{Abstract}\label{abstract}
\addcontentsline{toc}{subsection}{Abstract}

Indian packaged food labels do not share a common ingredient vocabulary.
The same substance appears under regional names, transliterations, INS
codes, and brand-specific terms --- sometimes across labels from the
same manufacturer. No reference layer exists that maps these expressions
to shared identities. This report documents the construction of a first
ingredient variant corpus from two sources: 896 directly sampled
products collected from verified Indian market listings, and English
ingredient strings from OpenFoodFacts filtered to rows with a traceable
Indian product name. Both sets were processed through a constrained
parsing pipeline --- one atomic operation per API call, temperature set
to 0, explicit failure rather than approximation when data was
unavailable. After combining both sources and iterative cleaning, the
corpus contains 2,291 unique ingredient variant strings. These variants
are not noise to eliminate. They are documentation of how ingredient
identity is expressed in practice across Indian commercial food labels.
The question of which variants refer to the same ingredient --- and by
what logic --- is addressed in EMF Model (A R 2026).

\section{The Question That Starts
Everything}\label{the-question-that-starts-everything}

A computer cannot tell you whether rice is healthier than Maggi. Not
because the comparison is philosophically difficult, but because the
infrastructure required to answer it does not exist.

To answer the question, the system needs to know what is in both
products. To know what is in them, it needs ingredient data. To use
ingredient data, it needs to know that ``maida'' and ``refined wheat
flour'' refer to the same thing --- and that ``palm oil'' and
``palmite'' do not, even though automated similarity measures would
score them close. To know that, it needs a stable reference layer that
maps the names labels actually use to the identities they actually mean.

That reference layer does not exist for Indian packaged food. This
report documents the first step toward building it: a collection of
ingredient variant strings extracted from commercial Indian food labels,
captured as they appear, without flattening the diversity that makes
them what they are.

\section{Why the Diversity Is Not the
Problem}\label{why-the-diversity-is-not-the-problem}

A label from a major Indian snack brand might read:

\begin{quote}
Seasoning Mix \{Iodised Salt, Chilli Powder (1.1\%), \#Spices \&
Condiments, Onion, Maltodextrin, Wheat Flour, Milk Solids, Black Salt,
Tomato Powder {[}Tomato Paste, Anticaking Agent (INS 551){]}, Refined
Sugar, Hydrolyzed Vegetable Protein, Acidity Regulators (INS 296, INS
330, INS 334), Garlic, Anticaking Agent (INS 551), Flavour Enhancers
(INS 627, INS 631)\} And Iodised Salt.
\end{quote}

This is not poorly formatted data. This is a brand communicating
ingredient relationships to consumers across India's 22 official
languages and hundreds of regional contexts, within the structure FSSAI
Labelling Rules 2020 require. The nested brackets encode functional
relationships: ``Acidity Regulators'' governs three INS codes as a
category. ``Tomato Powder'' contains both a base ingredient and an
additive. A Tamil-speaking consumer and a Hindi-speaking consumer both
need to read this label correctly. The formatting serves them.

The goal of this project is not to make that label simpler. It is to
build the layer underneath it that makes it machine-queryable ---
without asking ITC, or any other brand, to change a word.

A substrate is the layer that makes other things possible to build.
Concrete is a substrate: you do not live in concrete, you live in the
building the concrete made possible. The substrate does not care what
the building looks like. IFID --- Indian Food Informatics Data --- is
being built as that layer for ingredient identity. Tamil names stay
Tamil. INS codes stay in their FSSAI-specified format. The nested
bracket structure a brand uses to communicate to its consumers stays
exactly as designed. The substrate sits underneath and makes them
interoperable: queryable as the same ingredient when that is what you
need, distinguishable as different expressions when that matters.

Coordination without convergence. That is the specific goal.

\section{The Wall, and Who Is Already Working on
It}\label{the-wall-and-who-is-already-working-on-it}

Everyone who works with Indian packaged food data hits the same wall
from a different direction.

The nutritionist has fifty product samples and is spending half her time
cleaning label data before she can begin her actual analysis. The
e-commerce platform has the same ingredient listed seventeen different
ways across seventeen brands and cannot build a consistent product
catalogue. The compliance team is manually reconciling ingredient
declarations across FSSAI requirements, retailer formats, and export
documentation --- separately, every time. The researcher who could build
a tool to flag allergen risks has found there is no labelled dataset to
train on.

None of these people are doing it wrong. The wall is not their failure.
The wall is that no shared ingredient identity layer exists.

The most serious open effort to build one globally is OpenFoodFacts
(OFF). OFF has documented food products across dozens of countries
through crowdsourced and scraped contributions. The scale of that work
is significant and the intent is the same as this project's: make food
data open, structured, and usable. The gap in Indian product coverage
that this report documents is not a gap in OFF's effort. It is a direct
reflection of how fragmented and underdocumented the Indian packaged
food space actually is --- which is precisely what makes the problem
worth working on, and precisely what makes collaboration across efforts
like these necessary.

\section{Two Sources, One Problem}\label{two-sources-one-problem}

Building the ingredient vocabulary required two separate collection
strategies, for the same underlying reason: no single existing source
reliably answers whether a product is a current, shelf-available Indian
packaged food with a verifiable ingredient list.

\subsection{Why OFF Could Not Be the Only
Source}\label{why-off-could-not-be-the-only-source}

OFF contains thousands of English ingredient lists for products with
Indian brand names. Those lists are valuable. But the dataset structure
does not reliably distinguish a product currently on Indian supermarket
shelves from an imported variant, an export formulation, or a historical
listing no longer in distribution.

For the purpose of this corpus --- documenting what Indian consumers
actually encounter today --- that distinction matters. An ingredient
list attached to a product that is not in the Indian market does not
reflect the vocabulary Indian food systems use.

The null rates in the OFF data confirm the scale of the gap. Of 19,748
rows in the raw export:

\begin{itemize}
\tightlist
\item
  Only 4,104 pass a minimum filter: brand present, English product name
  present, English ingredient text present. That is 20.78 percent.
\item
  \texttt{ingredients\_text\_en} is the only ingredient column with
  coverage above 1 percent. All 29 other language columns combined add
  69 rows to that count.
\item
  6,905 rows have both a brand identifier and an English product name
  --- the product exists, it has a name --- but no ingredient text in
  any language. The gap is specifically at the ingredient field.
\item
  The four core macronutrient fields (energy, fat, protein,
  carbohydrates) have null rates between 65.61 and 66.00 percent across
  the full dataset.
\end{itemize}

{\def\LTcaptype{none} % do not increment counter
\begin{longtable}[]{@{}lrr@{}}
\toprule\noalign{}
Field & Non-null & Null \% \\
\midrule\noalign{}
\endhead
\bottomrule\noalign{}
\endlastfoot
energy\_value & 6,792 & 65.61 \\
fat\_value & 6,715 & 66.00 \\
proteins\_value & 6,737 & 65.89 \\
carbohydrates\_value & 6,759 & 65.77 \\
\end{longtable}
}

These numbers are not a criticism. They are a measurement of the space.
The gap in documented Indian food data is real, it is large, and it
exists because the underlying ecosystem is genuinely fragmented --- not
because anyone has failed to document it well enough.

What OFF does have --- the 4,104 rows with a brand, a product name, and
an English ingredient list --- is usable for this project. The product
name provides the minimum anchor needed to verify the product is Indian.
Those rows were taken through the same constrained parsing pipeline
described below, and their ingredient strings added to the vocabulary
set. The filter that kept them is described in the claims document.

\subsection{Why Direct Sampling Was
Necessary}\label{why-direct-sampling-was-necessary}

For a reliable picture of what is currently on Indian shelves, the
corpus needed to be collected directly. The methodology: select products
from companies with documented Indian market presence, retrieve
ingredient lists from verifiable online sources, extract and parse.

\textbf{Company and product selection.} Ten companies were selected
based on market presence across major packaged food categories ---
snacks, beverages, staples, dairy, condiments. Within each company,
selection moved through sub-brands (ITC's portfolio spans Aashirvaad,
Sunfeast, Bingo, YiPPee --- each with a different ingredient vocabulary)
and then to individual SKUs meeting four criteria:

\begin{enumerate}
\def\labelenumi{\arabic{enumi}.}
\tightlist
\item
  Ingredient list traceable on a whitelisted domain
\item
  Product available in the Indian market, not an export or international
  variant
\item
  Specificity to a single SKU, not a product range --- ``Aashirvaad
  Turmeric Powder 200g'' not ``Aashirvaad Spices''
\item
  One representative retained per formulation across pack sizes
\end{enumerate}

The third criterion produced the most rejections. References like
``Cadbury Chocolates'' or ``Aashirvaad Spices'' denote product families,
not individual items with specific ingredient lists. Every such
reference required disambiguation before it could enter the corpus.

After validation: 1,000 SKUs across 42 companies, 153 brands, 8
macro-categories.

\section{Why Standard Automated Parsing Fails
Here}\label{why-standard-automated-parsing-fails-here}

Before describing what the pipeline does, it is worth being precise
about why standard approaches do not work for this specific problem.

Palm oil and palmite are chemically and functionally distinct
ingredients. An automated similarity measure --- edit distance,
embedding cosine similarity, fuzzy matching --- would score them as
near-identical. Acting on that score would silently corrupt the
vocabulary.

Besan flour and chana dal are the same ingredient in different languages
and forms. A similarity measure that does not carry cultural and
linguistic knowledge would treat them as unrelated.

These are not edge cases. They are representative of how Indian food
labelling works: regional names, transliterations, preparation-state
variants, and INS codes all coexist on the same label, sometimes
referring to the same thing, sometimes to things that are genuinely
distinct. Standard clustering and normalisation algorithms cannot
reliably navigate this space. The cost of a silent error --- a wrong
merge, a missed distinction --- propagates forward into every analysis
built on the vocabulary.

The approach used here trades throughput for verifiability: one atomic
operation per API call, constrained to prevent approximation, with
explicit failure when the data is not there.

\section{The Extraction Pipeline}\label{the-extraction-pipeline}

\subsection{Constrained Retrieval}\label{constrained-retrieval}

The model retrieved ingredient lists from whitelisted domains only, in
priority order: brand official website, then Amazon India, BigBasket,
Blinkit. If no source returned the ingredient list, the output was
\texttt{DATA\_NOT\_FOUND}. The instruction was explicit: do not infer
typical ingredients from product category, do not approximate based on
similar products.

Temperature was set to 0. This means the model selects the
highest-probability token at each step and produces identical output for
identical input. The practical effect: if you run the same extraction
twice, you get the same result. Validation becomes tractable.
Fabrication through sampling variation is eliminated.

Batch size was tested across 1 to 20 SKUs per call. At batch sizes above
10, constraint violations increased measurably --- the model began
returning approximations instead of \texttt{DATA\_NOT\_FOUND} for
products it could not find, and formatting inconsistencies appeared. Six
SKUs per batch produced the best balance of throughput and constraint
adherence.

Results across 1,000 SKUs:

\begin{itemize}
\tightlist
\item
  First pass: 871 successful extractions (87.1\%), 129
  \texttt{DATA\_NOT\_FOUND}
\item
  Second pass on the 129 failures: 25 additional extractions, 104
  persistent failures
\item
  Final corpus: 896 extracted (89.6\%), 104 excluded
\end{itemize}

The 104 persistent failures validate that the constraint held. Those
products either had no verifiable online ingredient list or were no
longer in active distribution. The pipeline returned an explicit gap
rather than a filled approximation. An explicit gap can be addressed
later. A fabricated entry corrupts the vocabulary invisibly.

Manual audit of 90 extractions from the 896: 1 error identified (the
model merged content from two adjacent sections of a product page).
Error rate: 1 in 896 (0.11 percent).

\subsection{Structure-Preserving
Parsing}\label{structure-preserving-parsing}

The 896 extracted ingredient lists were not fed to the pipeline as whole
strings. Each list went through parsing as a discrete operation, because
the structure of Indian food labels encodes relationships that naive
splitting destroys.

Consider what happens when a comma-splitter treats every comma equally:

\textbf{Input:}
\texttt{Acidity\ Regulators\ (INS\ 296,\ INS\ 330,\ INS\ 334)}

\textbf{Naive output:}

\begin{itemize}
\tightlist
\item
  \texttt{Acidity\ Regulators\ (INS\ 296}
\item
  \texttt{INS\ 330}
\item
  \texttt{INS\ 334)}
\end{itemize}

The functional context --- that INS 296, 330, and 334 are all acidity
regulators --- is gone. The fragments \texttt{INS\ 330} and
\texttt{INS\ 334)} have no meaning without it.

The structure-aware parser tracks nesting depth. It splits on commas
only at depth zero --- the root level. Everything inside brackets is
treated as a unit until the brackets close. Applied to the same input:

\textbf{Structure-aware output:}
\texttt{Acidity\ Regulators\ (INS\ 296,\ INS\ 330,\ INS\ 334)} ---
intact, ready for decomposition with context preserved.

896 ingredient lists → 2,926 parsed strings with functional
relationships intact.

\subsection{Artifact Removal}\label{artifact-removal}

Each of the 2,926 strings went through a single-purpose cleaning pass:
remove presentation artifacts, preserve identity information.

Removed: percentage values (\texttt{55.7\%} --- quantity, not identity),
marketing text (\texttt{BINGO!}, \texttt{NEW!}), usage annotations
(\texttt{\#Used\ As\ Natural\ Flavouring\ Agent}).

Preserved: INS codes and E-numbers (regulatory identifiers), preparation
specifications (\texttt{Salt\ (Iodised)} --- the bracketed term
distinguishes a specific variety), functional classifications
(\texttt{Acidity\ Regulator}, \texttt{Emulsifier}).

The distinction matters because it is not always obvious.
\texttt{55.7\%} is presentation --- removing it loses nothing about what
the ingredient is. \texttt{(Iodised)} is identity --- removing it
collapses iodised salt and table salt into the same entry, which is
wrong.

Batch sizes for this stage were set inversely to string complexity:
short strings processed in batches of 40, complex multi-bracket strings
processed one at a time. Attention dilution at scale produces the same
constraint violations as in the retrieval stage.

\subsection{Semantic Decomposition}\label{semantic-decomposition}

After cleaning, compound structures were decomposed with context
propagation. Each atomic operation took one compound and returned its
components, with the functional classification carried forward to each:

\textbf{Input:} \texttt{Flavour\ Enhancers\ (INS\ 627,\ INS\ 631)}\\
\textbf{Output:} \texttt{Flavour\ Enhancer\ INS\ 627},
\texttt{Flavour\ Enhancer\ INS\ 631}

\textbf{Input:}
\texttt{Stabilizing\ \&\ Emulsifying\ Agents\ (412,\ 410,\ 407,\ 471,\ 466)}\\
\textbf{Output:} \texttt{Stabilizer\ INS\ 412},
\texttt{Stabilizer\ INS\ 410}, \texttt{Stabilizer\ INS\ 407},
\texttt{Emulsifier\ INS\ 471}, \texttt{Stabilizer\ INS\ 466}

\textbf{Input:}
\texttt{Black\ Pepper\ Powder,\ Ginger\ Powder,\ Clove\ Powder}\\
\textbf{Output:} unchanged --- already atomic

2,926 cleaned strings → 3,452 decomposed ingredient mentions → 1,987
unique variants after deduplication across all 896 products from the
sampling pipeline.

The full transformation for one product (Bingo Original Style, ITC Ltd.)
produced 21 ingredient mentions, including:

\begin{quote}
`Black Salt', `Chilli', `Citric Acid (INS 330)', `Disodium Guanylate
(INS 627)', `Disodium Inosinate (INS 631)', `Garlic', `Hydrolyzed
Vegetable Protein', `Maida', `Malic Acid (INS 296)', `Maltodextrin',
`Milk Solids', `Onion', `Palm Oil', `Potato', `Salt', `Silicon Dioxide
(INS 551)', `Spices and Condiments', `Sugar', `Tartaric Acid (INS 334)',
`Tomato'
\end{quote}

\section{Combining the Two Sources}\label{combining-the-two-sources}

The sampling pipeline produced 1,987 unique variant strings from 896
directly collected products. The OFF pipeline --- 4,104 rows filtered to
those with a verifiable product name, passed through the same
constrained parsing stages --- contributed an additional set of
ingredient strings from a different cross-section of the label space.

Combined and deduplicated across both sources, then cleaned through
multiple iterative audit rounds (documented in Appendix A), the final
corpus contains \textbf{2,291 unique ingredient variant strings}.

These are not errors to correct or synonyms to collapse. They are
documentation of how ingredient identity is expressed across Indian
commercial food labels. The same ingredient appears in multiple forms
because Indian food labelling reflects genuine linguistic and cultural
diversity:

\begin{itemize}
\tightlist
\item
  \texttt{Chilli} / \texttt{Chili} / \texttt{Chillies} --- orthographic
  variants, all in use
\item
  \texttt{Maida} / \texttt{Refined\ Wheat\ Flour} /
  \texttt{All-Purpose\ Flour} --- the same ingredient across language
  registers
\item
  \texttt{Onion\ Powder} / \texttt{Dried\ Onion} /
  \texttt{Dehydrated\ Onion} --- preparation-state variants
\item
  \texttt{INS\ 330} / \texttt{Citric\ Acid} /
  \texttt{Acidity\ Regulator\ INS\ 330} --- the same compound at
  different levels of regulatory specificity
\item
  \texttt{Iodised\ Salt} / \texttt{Salt\ (Iodised)} /
  \texttt{Table\ Salt} --- formatting alternatives for the same
  distinction
\end{itemize}

Each of these variants appears on labels that consumers read, regulators
review, and supply chains track. The infrastructure this project is
building needs to work with all of them --- not by picking one as
canonical and discarding the rest, but by organising them so that a
query for any one returns the right set.

The Tamil name on a label stays Tamil. The INS code stays in its FSSAI
format. The regional cultivar name stays as the brand printed it. The
substrate underneath makes them queryable as the same ingredient when
that is what the question requires.

\section{What This Corpus Makes
Possible}\label{what-this-corpus-makes-possible}

The output of this report is two open datasets:

\begin{itemize}
\tightlist
\item
  896 SKUs with verified ingredient lists, sourced from 42 companies
  across 8 macro-categories
\item
  Ingredient variant strings extracted from OFF data, filtered to rows
  with a verifiable Indian product name, cleaned through the same
  pipeline
\end{itemize}

Combined: a documented vocabulary of 2,291 unique ingredient expressions
from Indian packaged food labels, with extraction methodology,
constraint architecture, and quality validation documented in full.

The next question the corpus raises is: which of these 2,291 variants
refer to the same ingredient, and by what logic? \texttt{Maida} and
\texttt{Refined\ Wheat\ Flour} are the same substance.
\texttt{Palm\ Oil} and \texttt{Palmite} are not, despite surface
similarity. \texttt{Besan\ flour} and \texttt{chana\ dal} are related
but distinct in preparation state. Answering that question requires a
classification framework capable of handling identity, equivalence, and
subset relationships across a space where standard similarity measures
are unreliable.

That framework --- the EMF Model (Energy, Matter, Function) --- is
defined in A R (2026). Further progress on the mapping problem is
deferred to future reports.

\section{Claims and Verification}\label{claims-and-verification}

All numerical claims in this report are independently verifiable against
the source datasets. The full claims list with evidence per claim is
available at

\subsection{Claims}\label{claims}

{\def\LTcaptype{none} % do not increment counter
\begin{longtable}[]{@{}
  >{\raggedright\arraybackslash}p{(\linewidth - 2\tabcolsep) * \real{0.5000}}
  >{\raggedright\arraybackslash}p{(\linewidth - 2\tabcolsep) * \real{0.5000}}@{}}
\toprule\noalign{}
\begin{minipage}[b]{\linewidth}\raggedright
ID
\end{minipage} & \begin{minipage}[b]{\linewidth}\raggedright
Claim
\end{minipage} \\
\midrule\noalign{}
\endhead
\bottomrule\noalign{}
\endlastfoot
OFF.C.01 & Of 19,748 rows in the raw OpenFoodFacts export, 4,104 pass
the minimum filter (brand, product name in English, ingredient text in
English), a pass rate of 20.78 percent. \\
OFF.C.02 & ingredients\_text\_en is the only ingredient column with
coverage above 1 percent. It has 4,592 non-null rows (23.25 percent).
All 29 other language columns combined add 69 additional rows. \\
OFF.C.03 & 6,905 rows have both a brand identifier and an English
product name but no ingredient text in any language. The data gap is at
the ingredient field, not at product identity. \\
OFF.C.04 & The four core macronutrient fields have null rates between
65.61 and 66.00 percent across all 19,748 rows: energy\_value 65.61
percent (6,792 non-null), fat\_value 66.00 percent (6,715 non-null),
proteins\_value 65.89 percent (6,737 non-null), carbohydrates\_value
65.77 percent (6,759 non-null). \\
OFF.C.05 & The three Hindi language columns have the following non-null
counts across 19,748 rows: product\_name\_hi 111, ingredients\_text\_hi
11, generic\_name\_hi 2. \\
OFF.C.06 & Replacing product\_name\_en OR generic\_name\_en with
product\_name\_en alone as a filter condition reduces the output from
4,105 rows to 4,104 rows. generic\_name\_en contributes one unique
row. \\
OFF.C.07 & The raw dataset has 486 columns. The filtered dataset retains
4 columns: product\_name\_en, brands, brands\_tags, and
ingredients\_text\_en. \\
SAMP.C.01 & The sampling corpus spans 42 companies, 153 consumer-facing
brands, and 896 SKUs across 8 product macro-categories and 30
sub-categories. \\
SAMP.C.02 & The five highest-SKU companies --- Tata Consumer Products
(104), Amul / GCMMF (82), Haldiram's (68), Hindustan Unilever (67), and
ITC Ltd (65) --- account for 386 SKUs, or 43.1 percent of the 896-SKU
corpus. \\
SAMP.C.03 & SKU distribution across eight macro-categories derived from
top-3 category fields per company: beverages (200), sweets and desserts
(176), staples and spices (174), ready to eat and ready to cook (100),
snacks and namkeen (67), pantry and condiments (47), health and wellness
(43), dairy and breakfast (30). These sum to 837 of 896 total SKUs; the
remaining 59 fall into sub-categories not captured in the top-3 field
per company. \\
SAMP.C.04 & Of 1,000 SKUs submitted for extraction, 871 returned
successful ingredient lists on first pass (87.1 percent). A second-pass
retry on the 129 failures yielded 25 additional extractions (2.5
percent). Final corpus: 896 successful extractions (89.6 percent). 104
SKUs returned DATA\_NOT\_FOUND across both passes and are excluded. \\
SAMP.C.05 & Manual audit of 90 extractions from the 896-SKU corpus
identified 1 hallucination instance. Rate: 1 in 896 (0.11 percent). \\
SAMP.C.06 & Four SKU validation criteria were applied before extraction:
(1) ingredient list traceable within whitelisted domains; (2) product
available in the Indian market, not an export or international variant;
(3) specificity to a single SKU, not a product range; (4) one
representative retained per formulation across pack sizes. \\
SAMP.C.07 & Extraction operated under five constraints: (1) temperature
= 0; (2) domain whitelist: brand official website, Amazon India,
BigBasket, Blinkit, in priority order; (3) DATA\_NOT\_FOUND returned
when all sources exhausted; (4) JSON output schema enforced with four
required fields (product\_name, ingredient\_list, source\_url,
confidence); (5) brand official website given precedence over retailer
listings on conflict. \\
SAMP.C.08 & Batch sizes from 1 to 20 SKUs per API call were tested. 6
SKUs per batch was identified as optimal. Batches exceeding 10 SKUs
produced measurably increased constraint violations including
inappropriate DATA\_NOT\_FOUND omissions and formatting
inconsistencies. \\
\end{longtable}
}

\subsection{Evidence Per Claim}\label{evidence-per-claim}

\subsubsection{OFF.C.01}\label{off.c.01}

Raw row count: 19,748. Filter applied: brands OR brands\_tags non-empty,
AND product\_name\_en non-empty, AND ingredients\_text\_en non-empty.
Rows passing all three conditions: 4,104. Pass rate: 20.78 percent. Rows
removed: 15,644 (79.22 percent).

\subsubsection{OFF.C.02}\label{off.c.02}

{\def\LTcaptype{none} % do not increment counter
\begin{longtable}[]{@{}
  >{\raggedright\arraybackslash}p{(\linewidth - 2\tabcolsep) * \real{0.4286}}
  >{\raggedleft\arraybackslash}p{(\linewidth - 2\tabcolsep) * \real{0.5714}}@{}}
\toprule\noalign{}
\begin{minipage}[b]{\linewidth}\raggedright
Column
\end{minipage} & \begin{minipage}[b]{\linewidth}\raggedleft
Non-null rows
\end{minipage} \\
\midrule\noalign{}
\endhead
\bottomrule\noalign{}
\endlastfoot
ingredients\_text\_en & 4,592 \\
ingredients\_text\_fr & 94 \\
ingredients\_text\_de & 15 \\
ingredients\_text\_hi & 11 \\
All remaining 26 language columns (de-duplicated against English) &
69 \\
\end{longtable}
}

Pooling all 30 ingredient language columns yields 4,661 rows with any
ingredient text, against 4,592 for English alone.

\subsubsection{OFF.C.03}\label{off.c.03}

Rows passing (brands OR brands\_tags) AND product\_name\_en: 11,009. Of
these, rows also passing ingredients\_text\_en: 4,104. Rows with brand
and name but no ingredient text: 6,905.

\subsubsection{OFF.C.04}\label{off.c.04}

{\def\LTcaptype{none} % do not increment counter
\begin{longtable}[]{@{}lrrr@{}}
\toprule\noalign{}
Field & Null & Non-null & Null \% \\
\midrule\noalign{}
\endhead
\bottomrule\noalign{}
\endlastfoot
energy\_value & 12,956 & 6,792 & 65.61 \\
fat\_value & 13,033 & 6,715 & 66.00 \\
proteins\_value & 13,011 & 6,737 & 65.89 \\
carbohydrates\_value & 12,989 & 6,759 & 65.77 \\
\end{longtable}
}

Computed on the full 19,748-row dataset.

\subsubsection{OFF.C.05}\label{off.c.05}

{\def\LTcaptype{none} % do not increment counter
\begin{longtable}[]{@{}lrrr@{}}
\toprule\noalign{}
Column & Non-null & Null & Null \% \\
\midrule\noalign{}
\endhead
\bottomrule\noalign{}
\endlastfoot
product\_name\_hi & 111 & 19,637 & 99.44 \\
ingredients\_text\_hi & 11 & 19,737 & 99.94 \\
generic\_name\_hi & 2 & 19,746 & 99.99 \\
\end{longtable}
}

Computed on the full 19,748-row dataset.

\subsubsection{OFF.C.06}\label{off.c.06}

{\def\LTcaptype{none} % do not increment counter
\begin{longtable}[]{@{}
  >{\raggedright\arraybackslash}p{(\linewidth - 4\tabcolsep) * \real{0.3000}}
  >{\raggedright\arraybackslash}p{(\linewidth - 4\tabcolsep) * \real{0.3000}}
  >{\raggedleft\arraybackslash}p{(\linewidth - 4\tabcolsep) * \real{0.4000}}@{}}
\toprule\noalign{}
\begin{minipage}[b]{\linewidth}\raggedright
Filter
\end{minipage} & \begin{minipage}[b]{\linewidth}\raggedright
Condition
\end{minipage} & \begin{minipage}[b]{\linewidth}\raggedleft
Result
\end{minipage} \\
\midrule\noalign{}
\endhead
\bottomrule\noalign{}
\endlastfoot
Filter A & (brands OR brands\_tags) AND (product\_name\_en OR
generic\_name\_en) AND ingredients\_text\_en & 4,105 rows \\
Filter B & (brands OR brands\_tags) AND product\_name\_en AND
ingredients\_text\_en & 4,104 rows \\
\end{longtable}
}

Difference: 1 row. That row had generic\_name\_en populated and
product\_name\_en empty. In the 4,105-row set, generic\_name\_en has 451
non-null values (10.99 percent non-null, 89.01 percent null).

\subsubsection{OFF.C.07}\label{off.c.07}

Raw column count: 486. Columns retained after filter: product\_name\_en,
brands, brands\_tags, ingredients\_text\_en. Column count in working
dataset: 4. The 482 removed columns include all non-English name and
ingredient variants, all nutrient sub-fields, environmental scores,
packaging fields, and contributor metadata.

\subsubsection{SAMP.C.01}\label{samp.c.01}

Roster file header: Total SKUs: 896 \textbar{} Brands: 153 \textbar{}
Companies: 42 \textbar{} Parent cats: 8 \textbar{} Sub-cats: 30.

\subsubsection{SAMP.C.02}\label{samp.c.02}

{\def\LTcaptype{none} % do not increment counter
\begin{longtable}[]{@{}lr@{}}
\toprule\noalign{}
Company & SKUs \\
\midrule\noalign{}
\endhead
\bottomrule\noalign{}
\endlastfoot
Tata Consumer Products & 104 \\
Amul / GCMMF & 82 \\
Haldiram's & 68 \\
Hindustan Unilever & 67 \\
ITC Ltd & 65 \\
\textbf{Total (top 5)} & \textbf{386} \\
\end{longtable}
}

386 / 896 = 43.1 percent of corpus.

\subsubsection{SAMP.C.03}\label{samp.c.03}

Summed from top-3 category fields across all 42 company rows in roster.
Sum: 837. Difference from 896: 59 SKUs assigned to sub-categories below
each company's top three.

\subsubsection{SAMP.C.04}\label{samp.c.04}

First pass: 871 extracted, 129 DATA\_NOT\_FOUND. Second pass on 129: 25
additional, 104 persistent DATA\_NOT\_FOUND. Total extracted: 896. Total
excluded: 104. Pass rate: 896 / 1000 = 89.6 percent.

\subsubsection{SAMP.C.05}\label{samp.c.05}

Audit sample: 90 SKUs. Errors found: 1 (model merged content from
multiple webpage sections). Rate: 1 / 896 = 0.0011.

\subsubsection{SAMP.C.06}\label{samp.c.06}

Four criteria applied at SKU selection stage. Documented rejection
categories: product range references requiring disambiguation to
individual SKU (e.g., ``Aashirvaad Spices'' to ``Aashirvaad Turmeric
Powder 200g''); products not in Indian market distribution.

\subsubsection{SAMP.C.07}\label{samp.c.07}

Five constraints applied uniformly to all 1,000 attempted SKUs. System
instruction for DATA\_NOT\_FOUND: ``If ingredient list not found in
whitelisted domains, return DATA\_NOT\_FOUND. DO NOT infer typical
ingredients from product category. DO NOT approximate based on similar
products.''

\subsubsection{SAMP.C.08}\label{samp.c.08}

Batch sizes 1--20 tested during extraction development. Optimal: 6 SKUs
per batch. Violations at \textgreater10 SKUs: DATA\_NOT\_FOUND omissions
and formatting inconsistencies.

\section*{Appendix A: Sample Cleaning
Rounds}\label{appendix-a-sample-cleaning-rounds}
\addcontentsline{toc}{section}{Appendix A: Sample Cleaning Rounds}

The iterative cleaning process that produced the final 2,291 variant set
from the combined corpus was not individually logged. Individual round
logs were not maintained by design: the changes involved --- correcting
a transliteration typo, removing a fragment like \texttt{dried-powder}
that parsed as an ingredient but was a formatting artifact, deciding
whether \texttt{monohydrate} belonged in the corpus at all --- were too
granular and numerous to document round by round without the log itself
becoming unmanageable.

What follows is a representative excerpt from the audit scripts used
during this process. It shows what the review actually looked like:
automated flagging, human decision at each boundary case, iterative
convergence toward a clean set.

\textbf{One audit pass --- executive summary:}

\begin{verbatim}
=============================================
      AI AUDIT EXECUTIVE SUMMARY
=============================================
Total Entries Audited : 709

APPROVED               :  623 (87.9%)
MODIFIED               :   51 (7.2%)
INVALID                :   35 (4.9%)
=============================================
\end{verbatim}

\textbf{Entries flagged as INVALID --- strings that were not
ingredients:}

\begin{verbatim}
atlantic · center-filling · cfu · chips · compound · dessert
dried-powder · dry · energy · flakes · food-additives · lubrication
moisture · monohydrate · mononitrate · only · pizza · plant-base
powder-mix · preservative · protein · savouries · test · toppings
vegetable · vegetable-mix · ...
\end{verbatim}

\textbf{Interactive kill review --- the monohydrate decision:}

The boundary cases required a human in the loop. \texttt{monohydrate}
and \texttt{mononitrate} are not ingredients --- they are suffixes that
appear on ingredient labels (as in \texttt{thiamine\ mononitrate}) but
carry no identity when extracted alone. They were saved on first pass,
then removed on second review:

\begin{verbatim}
KILL 'monohydrate'? (y/n): n
Saving 'monohydrate'...

 Surgery Complete. Your files are now 'Steel'. 

[second pass]

KILL 'monohydrate'? (y/n): y
Executing 'monohydrate'...

 Surgery Complete. Your files are now 'Steel'. 
\end{verbatim}

\textbf{Reclassification pass --- where the judgments were not
straightforward:}

\begin{verbatim}
ITEM: fish
FROM: Additives & Functional  →  TO: Proteins & Meats
Accept? y 

ITEM: gluten
FROM: Staples (Grains/Dals)  →  TO: Proteins & Meats
Accept? n   Added to manual review.

ITEM: spirulina
FROM: Additives & Functional  →  TO: Fruits, Veg & Botanicals
Accept? n   Added to manual review.

ITEM: fava-bean-protein
FROM: Additives & Functional  →  TO: Proteins & Meats
Accept? n   Added to manual review.
\end{verbatim}

\texttt{gluten}, \texttt{spirulina}, and \texttt{fava-bean-protein} are
examples where the automated reclassification suggestion was defensible
but not settled --- each sits at a category boundary that requires a
classification framework to resolve, not a cleaning pass. They were held
for the mapping stage.

\textbf{Final state after all cleaning rounds:}

\begin{verbatim}
==================================================
     FINAL MONOGRAPH DATA SUMMARY
==================================================
Total Raw Variants (TSV)    : 46,635
Unique Canonical Units      : 662
==================================================
\end{verbatim}

The 46,635 raw variants and 662 canonical units shown here are from the
OFF monograph specifically --- a separate but parallel cleaning process
applied to the OFF-derived strings. The 2,291 figure reported in the
main body is the combined and deduplicated variant set from both
sources, prior to canonical mapping. These are different counts at
different stages of the pipeline and are not in conflict.

\section*{Acknowledgements}\label{acknowledgements}
\addcontentsline{toc}{section}{Acknowledgements}

My deepest gratitude to Mr.~Krishna, whose constancy forms the
foundation upon which all my work, including this, quietly rests.
Salutations to the Goddess who dwells in all beings in the form of
intelligence. I bow to her again and again.

We are deeply grateful to all contributors of OFF Dataset - one of the
core sources which our efforts build upon. Thank you for all that you
do. This report was prepared as part of the Indian Food Informatics Data
(IFID) project at the Interdisciplinary Systems Research Lab (ISRL).

\section*{References}\label{references}
\addcontentsline{toc}{section}{References}

\protect\phantomsection\label{refs}
\begin{CSLReferences}{1}{1}
\bibitem[\citeproctext]{ref-arIdentityTransformationFunction2026}
A R, Lalitha. 2026. \emph{Identity, {Transformation}, and {Function}: {A
Tri-Axial Model} for the {Classification} of {Food Ingredient
Identity}}. Zenodo. \url{https://doi.org/10.5281/zenodo.18714527}.

\end{CSLReferences}




\end{document}
